\documentclass[11pt]{article}
\usepackage[margin=1in]{geometry}
\usepackage{amsmath,amssymb,amsthm}
\usepackage{graphicx}
\usepackage[hidelinks]{hyperref}
\setlength{\emergencystretch}{2em}

\newtheorem{theorem}{Theorem}
\newtheorem{lemma}{Lemma}
\newtheorem{corollary}{Corollary}
\newtheorem{remark}{Remark}

\newcommand{\C}{\mathbb C}
\newcommand{\dist}{\operatorname{dist}}
\newcommand{\otmin}{\otimes_{\min}}
\newcommand{\Prim}{\operatorname{Prim}}

\title{Uncountable Tensor Products and Countable Capture}
\author{Run \texttt{fa\_banach\_001}, lane 19}
\date{June 17, 2026}

\begin{document}
\maketitle

\section*{Source Problem}

The source paper is Jinghao Huang, Karimbergen Kudaybergenov, and Rui Liu,
\emph{The ball-covering property of von Neumann algebras and noncommutative
symmetric spaces}, arXiv:2311.14257.  After proving that a commutative unital
\(C^*\)-algebra \(C(K)\) has BCP/UBCP exactly when \(K\) has a countable
\(\pi\)-basis, the authors ask for a noncommutative version:

\begin{center}
\includegraphics[width=0.92\textwidth]{figures/open_problem_crop.png}
\end{center}

\section*{Statement}

\begin{theorem}
\label{thm:tensor-capture}
Let \((A_\gamma)_{\gamma\in\Gamma}\) be a family of unital \(C^*\)-algebras.
Assume that uncountably many \(A_\gamma\)'s are nontrivial.  Let
\[
        A=\bigotimes_{\gamma\in\Gamma}^{\min} A_\gamma
\]
be the unital minimal infinite tensor product with respect to the unit
embeddings.  Then \(A\) has the countable-capture property.  In particular,
\(A\) fails BCP.
\end{theorem}

\begin{corollary}
\label{cor:uhf}
For every uncountable set \(\Gamma\), the unital \(C^*\)-algebra
\[
        \bigotimes_{\gamma\in\Gamma}^{\min} M_2
\]
has countable capture and fails BCP.  This algebra is simple, so its primitive
ideal space is a point and has a countable \(\pi\)-base.
\end{corollary}

\section*{Idea of the Proof}

Every element of a unital infinite tensor product has countable support: it is
a norm limit of finite-support tensors, so the union of the approximating
finite supports is countable.  Thus a countable family of proposed centres
lives inside a countable subproduct.

Choose an unused nontrivial factor \(A_\gamma\).  Inside it choose a
self-adjoint contraction \(v\) whose spectrum contains both \(-1\) and \(1\).
Then there are states \(\rho_+\) and \(\rho_-\) with
\[
        \rho_+(v)=1,\qquad \rho_-(v)=-1.
\]
For a centre \(b\) in the old subproduct and any scalar \(\lambda\), the
slice maps \(\operatorname{id}\otimes\rho_\pm\) give
\[
        \|b-\lambda(1\otimes v)\|
        \ge
        \max\{\|b-\lambda 1\|,\|b+\lambda 1\|\}
        \ge
        \|b\|.
\]
The last inequality is just the triangle inequality applied to
\((b-\lambda1)+(b+\lambda1)=2b\).  Therefore
\[
        \dist(b,\C(1\otimes v))=\|b\|
\]
for every centre \(b\), which is exactly the quotient-line characterization
of countable capture from the previous packet.

\section*{Preliminaries}

We use the countable-capture characterization proved in the packet
\[
\texttt{2311.14257\_countable\_capture\_characterization}.
\]
For a nonzero normed space \(X\), countable capture is equivalent to the
following condition: for every countable \(C\subset X\), there exists
\(u\in S_X\) such that
\[
        \dist(c,\C u)=\|c\|
        \qquad(c\in C).
\]

\begin{lemma}[countable support]
\label{lem:support}
Let
\[
        A=\bigotimes_{\gamma\in\Gamma}^{\min} A_\gamma
\]
be a unital infinite tensor product.  For every \(a\in A\), there is a
countable subset \(\Lambda\subseteq\Gamma\) such that
\[
        a\in A_\Lambda
        :=
        \bigotimes_{\gamma\in\Lambda}^{\min} A_\gamma
\]
under the canonical unit embedding \(A_\Lambda\hookrightarrow A\).
Consequently every countable subset of \(A\) is contained in \(A_\Lambda\)
for some countable \(\Lambda\subseteq\Gamma\).
\end{lemma}

\begin{proof}
The infinite tensor product is the norm closure of the algebraic union of the
finite subproducts \(A_F\), \(F\subseteq\Gamma\) finite.  Given \(a\in A\),
choose finite-support tensors \(a_n\) with \(a_n\to a\).  Let
\[
        \Lambda=\bigcup_{n=1}^\infty \operatorname{supp}(a_n).
\]
Then \(\Lambda\) is countable, every \(a_n\) belongs to \(A_\Lambda\), and
\(A_\Lambda\) is closed in \(A\).  Hence \(a\in A_\Lambda\).  Applying this
to countably many elements and taking the union of the resulting countable
sets gives the final assertion.
\end{proof}

\begin{lemma}[two-state element in a nontrivial unital algebra]
\label{lem:two-state}
Let \(B\) be a nontrivial unital \(C^*\)-algebra.  Then there is a
self-adjoint element \(v\in B\) with \(\|v\|=1\), and states
\(\rho_+,\rho_-\) on \(B\), such that
\[
        \rho_+(v)=1,\qquad \rho_-(v)=-1.
\]
\end{lemma}

\begin{proof}
Choose a self-adjoint non-scalar element \(h\in B\).  Its spectrum contains
at least two points.  Let
\[
        m=\min\sigma(h),\qquad M=\max\sigma(h),
\]
so \(m<M\), and define
\[
        v=\frac{2h-(M+m)1}{M-m}.
\]
Then \(v=v^*\), \(\|v\|=1\), and \(-1,1\in\sigma(v)\).  The evaluation
characters at \(-1\) and \(1\) on \(C^*(v)\cong C(\sigma(v))\) are states;
extend them to states \(\rho_-\) and \(\rho_+\) on \(B\).
\end{proof}

\begin{lemma}[fresh tensor factor distance]
\label{lem:distance}
Let \(B\) and \(D\) be unital \(C^*\)-algebras, let \(v\in D\) and states
\(\rho_\pm\) satisfy Lemma~\ref{lem:two-state}.  For every \(b\in B\) and
every \(\lambda\in\C\),
\[
        \|b\otimes 1-\lambda(1\otimes v)\|_{B\otmin D}
        \ge
        \|b\|.
\]
Consequently, for \(u=1\otimes v\),
\[
        \dist(b\otimes1,\C u)=\|b\|.
\]
\end{lemma}

\begin{proof}
The slice maps \(\operatorname{id}_B\otimes\rho_\pm\) are contractive.  Hence
\[
\begin{aligned}
        \|b\otimes 1-\lambda(1\otimes v)\|
        &\ge
        \max\{\|b-\lambda\rho_+(v)1\|,
              \|b-\lambda\rho_-(v)1\|\}\\
        &=
        \max\{\|b-\lambda1\|,\|b+\lambda1\|\}.
\end{aligned}
\]
By the triangle inequality,
\[
        2\|b\|
        =
        \|(b-\lambda1)+(b+\lambda1)\|
        \le
        \|b-\lambda1\|+\|b+\lambda1\|
        \le
        2\max\{\|b-\lambda1\|,\|b+\lambda1\|\}.
\]
Thus the displayed norm is at least \(\|b\|\) for every \(\lambda\).  Taking
the infimum over \(\lambda\) gives
\[
        \dist(b\otimes1,\C u)\ge \|b\|.
\]
The reverse inequality follows by taking \(\lambda=0\), and \(\|u\|=1\).
\end{proof}

\section*{Proof of the Main Theorem}

\begin{proof}[Proof of Theorem~\ref{thm:tensor-capture}]
Let \(C\subset A\) be countable.  By Lemma~\ref{lem:support}, there is a
countable \(\Lambda\subseteq\Gamma\) such that
\[
        C\subseteq A_\Lambda.
\]
Since uncountably many factors are nontrivial, choose
\(\gamma_0\in\Gamma\setminus\Lambda\) with \(A_{\gamma_0}\) nontrivial.
Apply Lemma~\ref{lem:two-state} in \(A_{\gamma_0}\), obtaining
\(v\in A_{\gamma_0}\) and states \(\rho_\pm\) with \(\rho_\pm(v)=\pm1\).

Let \(u\in A\) be the tensor which is \(v\) in the coordinate \(\gamma_0\)
and the unit in every other coordinate.  Then \(\|u\|=1\).  For each
\(c\in C\), view \(c\) as an element of
\[
        A_\Lambda\otmin A_{\gamma_0}
\]
with trivial \(\gamma_0\)-coordinate.  Lemma~\ref{lem:distance} gives
\[
        \dist(c,\C u)=\|c\|.
\]
The countable-capture characterization now implies that \(A\) has countable
capture.  The previous countable-capture obstruction packet then implies
that \(A\) fails BCP.
\end{proof}

\begin{proof}[Proof of Corollary~\ref{cor:uhf}]
Apply Theorem~\ref{thm:tensor-capture} to \(A_\gamma=M_2\).  The resulting
algebra is the uncountable CAR/UHF tensor product.  It is the direct limit of
the finite matrix algebras \(\bigotimes_{\gamma\in F}M_2\), over finite
\(F\subseteq\Gamma\), with unital full connecting maps; hence it is simple.
Therefore its primitive ideal space is a single point, which has a countable
\(\pi\)-base.
\end{proof}

\section*{Scope and Remaining Wall}

This theorem supplies a new internal-size obstruction in the countable
primitive-\(\pi\)-base region.  It is independent of primitive topology:
the uncountable \(M_2\)-tensor product is simple, yet fails BCP by countable
capture.

The result does not classify all simple nonseparable \(C^*\)-algebras.  It
applies when the algebra has a literal uncountable tensor-coordinate
structure, or when a future argument can reduce a countable centre family to
a separable subalgebra with a fresh two-state unit available in a commuting
or tensor-independent direction.

\section*{Novelty and Search Bounds}

The local run indexes were searched for tensor products, infinite tensor
products, UHF algebras, countable support, and arXiv:2311.14257.  Web searches
on June 17, 2026 for BCP with infinite tensor products, UHF algebras, and
uncountable tensor products found no exact statement of this theorem.  The
argument is an agent-produced application of the countable-capture
characterization.

\section*{Human Review Focus}

Reviewers should check:
\begin{enumerate}
    \item the countable-support lemma for unital infinite tensor products;
    \item the construction of \(v\) and the two endpoint states in a
          nontrivial unital \(C^*\)-algebra;
    \item the slice-map distance inequality;
    \item the simplicity assertion in the \(M_2\)-tensor corollary.
\end{enumerate}

\begin{thebibliography}{9}

\bibitem{HKL}
Jinghao Huang, Karimbergen Kudaybergenov, and Rui Liu,
\emph{The ball-covering property of von Neumann algebras and noncommutative
symmetric spaces},
arXiv:2311.14257, 2023.

\bibitem{Murphy}
Gerard J. Murphy,
\emph{C*-Algebras and Operator Theory},
Academic Press, 1990.

\bibitem{lane19capture}
Run \texttt{fa\_banach\_001}, lane 19,
\emph{When the Countable-Capture BCP Obstruction Exists},
solution packet
\texttt{2311.14257\_countable\_capture\_characterization}, 2026.

\bibitem{lane19obstruction}
Run \texttt{fa\_banach\_001}, lane 19,
\emph{A Countable-Capture Obstruction for the Ball-Covering Property},
solution packet
\texttt{2311.14257\_countable\_capture\_bcp\_obstruction}, 2026.

\end{thebibliography}

\end{document}
